\documentclass{article}

% ICLR 2027 submission style.
% Run `make template` once to download the official ICLR 2027 style files.
\usepackage{iclr2027_conference,times}

% Minimal math-command file for the Hidden Policy draft.
% Kept intentionally small so notation remains easy to follow.
\newcommand{\E}{\mathbb{E}}
\newcommand{\R}{\mathbb{R}}


\usepackage{hyperref}
\usepackage{url}
\usepackage{booktabs}
\usepackage{amsmath,amssymb,amsthm,mathtools,centernot}
\usepackage{microtype}
\usepackage{xspace}

\newtheorem{definition}{Definition}
\newtheorem{proposition}{Proposition}
\newcommand{\HP}{\textsc{Hidden Policy}\xspace}
\newcommand{\Rint}{\mathcal{R}}
\newcommand{\Rreveal}{\mathcal{R}^{\prime}}
\newcommand{\Hist}{\mathcal{H}}
\newcommand{\Act}{\mathcal{A}}
\newcommand{\Traj}{\mathcal{T}}

\title{Hidden Policies: The Risk Behind Risks}

% Anonymous submission. Replace after the review period if needed.
\author{Anonymous Authors}

\begin{document}
\maketitle

\begin{abstract}
Frontier models' safety and security are usually assessed through \textbf{what a model can do} and \textbf{what it has done}. Deployment decisions, however, depend on a tougher question: \textbf{what will the model do?}  We argue that answering this question requires understanding the policy structure that governs behaviors across conditions.
Specifically, we define a \emph{hidden policy} as a generalizing, condition-gated decision rule whose influence on action selection is behaviorally masked within an interaction regime, rendering it observationally equivalent to a different policy structure yet systematically distinguishable under a reachable condition or intervention.
%Specifically, we focus on \textit{hidden policy} which refers to a generalizing, condition-gated decision rule whose influence on action selection is behaviorally masked within an interaction regime, rendering it observationally equivalent to a public policy yet systematically distinguishable under an unknown condition or intervention.
This definition distinguishes hidden policies from previous backdoor attacks, hidden objectives, and misaligned behaviors, while providing a common mechanism for sleeper agents, sandbagging, alignment faking, scheming, and the conditional deployment of capabilities related to model preparedness.
However, this risk exposes a policy-identifiability challenge: \textbf{behavioral evidence can only reveal the underlying policy up to an observational equivalence class and therefore fails to predict future behaviors when supervision, incentives, access, or authentication change in deployment}. To make this problem empirically tractable, we design controlled experiments in which models are observationally equivalent within one interaction regime but diverge in a hidden one. Constructing cases that range from memorized responses and conditional associations to generalized policy structures creates testbeds that enable diagnosis of whether a hidden policy has formed.
% Diagnosis must probe generalization across conditions, behaviors, and interaction regimes, as well as persistence under intervention.
The same testbed also makes it possible to evaluate the removal of hidden policy: whether an intervention eliminates the underlying policy structure or merely invalidates a trigger, blocks the behavior, or suppresses its expression. This construct-diagnose-remove framework shifts safety and security evaluation from measuring past behavior alone toward identifying and controlling the policies that govern future behaviors.
\end{abstract}

\section{Introduction}

Frontier large language models (LLMs) can possess emerging capabilities without expressing them, and they can produce safe responses under evaluation without doing so when conditions change. These discrepancies become increasingly important as models become more situationally aware, more capable of long-horizon reasoning, and more deeply embedded in environments where supervision, tool access, incentives, and authentication vary. Existing safety or security studies offer striking examples of condition-dependent behavior: sleeper agents that switch behavior under deployment cues \citep{hubinger2024sleeper}, models that selectively comply during training to preserve behavior outside training \citep{greenblatt2024alignment}, models optimized around hidden objectives \citep{marks2025hidden}, alongside preparedness frameworks that track severe risks spanning biological and chemical capabilities, cybersecurity, automated AI research, long-horizon autonomy, safeguard circumvention, and loss of control \citep{openai2025preparedness,deepmind2025frontier,anthropic2025rsp}. Yet these studies provide limited evidence about whether models would independently initiate and carry out the corresponding harms, because the observed behaviors are typically elicited under certain conditions that researchers explicitly specify or implicitly embed in training and evaluation setups. This leaves a gap between capability and observed behavior: \textbf{Does possessing a dangerous capability necessarily mean that a model will engage in harmful action?}

To close this gap, we distinguish three questions: \textbf{(Q1)}  \emph{What can a model do?} concerns capability; \textbf{(Q2)} \emph{What has it done?} concerns behavior observed under particular conditions;  \textbf{(Q3)} and \emph{What will it do as conditions change?} concerns the policy governing future action.  Existing safety or security evaluations often address the first two while leaving the third underdefined: \textbf{behavioral evidence may only show what a model can do and what it did under specific conditions, but cannot identify a policy that governs what it will do under varying conditions.} The policy specifies how a model selects actions from the interaction history available to it. Formally, a model induces a policy
\begin{equation}
    \pi(a\mid h),
\end{equation}
where $h\in\Hist$ denotes the interaction history available to the model and $a\in\Act$ denotes its next action. Building on this policy-level view, we define \emph{hidden policy} as \textbf{a special }\textbf{policy whose influence on action selection is behaviorally masked within the interactions history available to an observer (the owner or regulator of the model), yet becomes distinguishable under certain unknown conditions or interventions.} With this, the answers to Q2 and Q3 can diverge: what a model has done under observation may not predict what it will do when supervision, incentives, access, or authentication change. This divergence, rather than dangerous capabilities, should be treated as the meta risk in frontier risk assessment and preparedness evaluation, as it is the underlying mechanism of them: a dangerous capability can engage in harmful action when a hidden policy targets to express it. Moreover, such a policy can emerge from the model's own learning dynamics or be deliberately implanted by humans through pre-training or fine-tuning, which complicates the issue itself.

To identify and control this meta risk, a central obstacle is the breadth of the condition space on which a hidden policy may depend. Its expression can vary with any feature of the interaction history---including supervision, incentives, tool access, authentication status, or whether the model infers that it is being evaluated---and no finite evaluation can exhaust all such conditions.  Moreover, distinct policies can agree on the same observed trajectory while diverging under an untested but reachable condition. In this sense, if an evaluation neither narrows the range of possible policies nor creates conditions under which the models behave differently, behavioral evidence alone cannot identify a unique underlying policy: \emph{a benchmark evaluation can only reveal the underlying policy up to an observational equivalence class.}

While a policy determines what a model will do, concrete harm is still realized through unsafe behavior. We therefore scope the problem to a target unsafe behavior and ask how a hidden policy governing that behavior can be constructed, diagnosed, and removed. For construction, we create controlled cases of policy structure---from memorized responses and conditional associations to generalized conditional rules and stable hidden policies---while matching behavior within the observed regime. Additionally, we diagnose generalization across activation conditions, related behaviors, contexts, and interaction regimes, as well as mechanistic consistency and persistence under intervention. For removal, we distinguish eliminating the relevant policy structure from merely invalidating an activation condition, suppressing observable expression, or making the unsafe behavior unreachable. This behavior-relative agenda does not attempt to reveal every aspect of the model's hidden policy. Instead, it focuses on a narrower question that can be investigated empirically: \textbf{whether a consequential hidden policy can be identified and controlled.}

Our contributions are summarized as follows. \textbf{(1) What is a hidden policy, and what does the observational equivalence capture?} We formalize a hidden policy regarding an observer's interaction regime and show that behavioral evidence can reveal the underlying policy only up to an observational equivalence class. \textbf{(2) Why are hidden policies important?} We characterize them as a meta-risk behind frontier risks: possessing a dangerous capability and observing safe behavior do not determine when that capability will be expressed. \textbf{(3) How can a hidden policy be constructed?} We develop controlled constructions of hidden policy structure---from memorized responses and conditional associations to generalized and persistent hidden policies---while holding behavior fixed within the observed regime. \textbf{(4) When a hidden policy is removed?} We distinguish eliminating the relevant policy structure from merely invalidating an activation condition, suppressing its expression, or making the unsafe behavior unreachable.

\section{What Is a Hidden Policy?}
\label{sec:definition}

%%%%%%%%%%%%%%%%%%%%%%%%%%%%%%%%%%%%%%%%%%%%%%%%%%%%%%%%%%%%%%%%%%%%%
%%%%%%%%%%%%%%%%%%%%%%%%%%%%%%%%%%%%%%%%%%%%%%%%%%%%%%%%%%%%%%%%%%%%%
%%%%%%%%%%%%%%%%%%%%%%%%%%%%%%%%%%%%%Capability, policy, and behavior
%%%%%%%%%%%%%%%%%%%%%%%%%%%%%%%%%%%%%%%%%%%%%%%%%%%%%%%%%%%%%%%%%%%%%
%%%%%%%%%%%%%%%%%%%%%%%%%%%%%%%%%%%%%%%%%%%%%%%%%%%%%%%%%%%%%%%%%%%%%

\subsection{Capability, policy, and behavior}

A central distinction in this work is between \emph{what a model can do}, \emph{what it actually does}, and \emph{how it decides what to do}. We refer to these as \emph{capability}, \emph{behavior}, and \emph{policy}, respectively. Although closely related, they describe different properties of a model.

\paragraph{Capability: what can the model do?}
A capability describes a class of behaviors that a model can reliably produce under suitable conditions. For example, a model may be capable of exploiting a software vulnerability if it can reliably perform the exploit when provided with the necessary information, tools, and opportunity. More generally, for a target behavior class $\mathcal{B}$ and a family of admissible elicitation conditions $\mathcal{E}$, capability asks whether there exists some $r\in\mathcal{E}$ under which the model can reliably produce a target behavior $b\in\mathcal{B}$.
% Capability therefore characterizes what behaviors are available to the model under suitable conditions, rather than which behavior it will select in a particular interaction.

\paragraph{Behavior: what did the model actually do?}
Behavior describes what is observed when the model interacts with its environment under a particular condition. We represent an interaction by a trajectory $\tau\in\Traj$, which records the sequence of observations and actions that occurs during the interaction.
% An observed behavior may be a single trajectory or a statistic computed from repeated trajectories under the same condition.
Behavior is therefore conditional on the circumstances in which the model is observed. The same model may exhibit different behaviors when, for example, supervision, incentives, access to tools, authentication, or other aspects of the interaction change.
% Observing behavior under one set of conditions tells us what the model did there, but does not by itself specify what it would do under another.

\paragraph{Policy: how the model decides what to do?}
To describe how behavior varies across conditions, we introduce the model's \emph{policy}. We represent an interaction as a sequence of observations and actions,
\begin{equation}
\tau=(o_1,a_1,o_2,a_2,\ldots,o_T,a_T)\in\Traj,
\end{equation}
where $o_t$ denotes the information observed by the model at step $t$ and $a_t\in\Act$ denotes its action. A policy $\pi\in\Pi$ is the decision rule that governs how the model selects actions throughout such an interaction. It may condition its decisions on the information observed so far, including user inputs, previous interactions, supervision, available tools, authentication or access state, and other information from which the model may infer its current situation. The policy therefore governs how the model acts across possible interaction conditions, rather than describing any single realized behavior.

We use an \emph{interaction regime} $r\subseteq\Traj$ to denote a class of interactions sharing a condition of interest. A regime may, for example, distinguish evaluation from deployment, monitored from unmonitored interactions, or authenticated from unauthenticated access. More generally, the choice of regime depends on the conditions relevant to the behavior being studied. We write $\Rint$ for the collection of interaction regimes under consideration. A trajectory $\tau$ represents a single realized interaction, whereas a regime $r$ identifies a class of such interactions. For a policy $\pi$ and regime $r\in\Rint$, we write
\begin{equation}
P_{\pi}^{r},
\qquad
\tau\sim P_{\pi}^{r},
\end{equation}
where $P_{\pi}^{r}$ denotes the distribution over trajectories induced by policy $\pi$ within regime $r$, and $\tau\sim P_{\pi}^{r}$ denotes a realized interaction drawn from that distribution. Thus, what we observe is generated by $(\pi,r)\longmapsto P_{\pi}^{r}$, but observing $P_{\pi}^{r}$ within a limited set of regimes can not identify a unique $\pi$. Different policies can be behaviorally indistinguishable where we observe them, yet diverge elsewhere. This ambiguity is the basis of observational equivalence and, ultimately, hidden policies.

%%%%%%%%%%%%%%%%%%%%%%%%%%%%%%%%%%%%%%%%%%%%%%%%%%%%%%%%%%%%%%%%%%%%%
%%%%%%%%%%%%%%%%%%%%%%%%%%%%%%%%%%%%%%%%%%%%%%%%%%%%%%%%%%%%%%%%%%%%%
%%%%%%%%%%%%%%%%%%%%%%%%%%%%%%%%%%%%%%%%%%%%%%%%%%%%%%%%%%%%%%%%%%%%%
%%%%%%%%%%%%%%%%%%%%%%%%%%%%%%%%%%%%%%%%%%%%%%%%%%%%%%%%%%%%%%%%%%%%%

\subsection{Observational equivalence and policy identification}
\label{sec}

An observer (e.g., the owner or regulator of the model) can examine a model only within a set of interaction regimes. Let $\mathcal{R}_{\mathrm{obs}}\subseteq\Rint$ denote the regimes available to the observer.

\begin{definition}[Observational equivalence]
Two policies $\pi_1,\pi_2\in\Pi$ are \emph{observationally equivalent} under $\mathcal{R}_{\mathrm{obs}}$, written
$\pi_1\sim{\mathcal{R}_{\mathrm{obs}}}\pi_2$, if
\begin{equation}
P_{\pi_1}^{r}=P_{\pi_2}^{r}
\qquad
\text{for every } r\in\mathcal{R}_{\mathrm{obs}}.
\end{equation}
\end{definition}

Observational equivalence partitions the candidate policy class into equivalence classes
\begin{equation}
[\pi]{\mathcal{R}_{\mathrm{obs}}}
:=
{\pi'\in\Pi :
\pi'\sim_{\mathcal{R}_{\mathrm{obs}}}\pi}.
\end{equation}

Crucially, two policies may be indistinguishable throughout the regimes available to the observer while behaving differently in another reachable regime:
\begin{equation}
\pi_1\sim_{\mathcal{R}{\mathrm{obs}}}\pi_2,
\qquad
P_{\pi_1}^{r^\star}\neq P_{\pi_2}^{r^\star}
\quad
\text{for some } r^\star\in\Rint\setminus\mathcal{R}_{\mathrm{obs}}.
\end{equation}

%%%%%%%%%%%%%%%%%%%%%%%%%%%%%%%%%%%%%%%%%%%%%%%%%%%%%%%%%%%%%%%%%%%%%
%%%%%%%%%%%%%%%%%%%%%%%%%%%%%%%%%%%%%%%%%%%%%%%%%%%%%%%%%%%%%%%%%%%%%
%%%%%%%%%%%%%%%%%%%%%%%%%%%%%%%%%%%%%%%%%%%%%%%%%%%%%%%%%%%%%%%%%%%%%
%%%%%%%%%%%%%%%%%%%%%%%%%%%%%%%%%%%%%%%%%%%%%%%%%%%%%%%%%%%%%%%%%%%%%

% \paragraph{1. Complete-policy comparison.}
% Both $\pi_H$ and $\pi_{\mathrm{ref}}$ are complete mappings from histories to action distributions. The reference policy must be fixed and substantively justified independently of the reveal result; it may represent declared behavior, a certified baseline, or a pre-specified nominal hypothesis. It is not a second component assumed to be stored inside the model. What is hidden is which complete policy the model implements. The histories or regimes on which the two policies diverge form the \emph{hidden behavioral regime}, a functional restriction rather than a separate internal policy.

% \paragraph{2. Reachable, systematic separation.}
% The reveal family must contain conditions available in deployment or through admissible interventions, rather than logically possible but operationally irrelevant counterfactuals. The measure $\mu$ declares which variations count, and $\alpha$ requires divergence over a non-negligible family rather than at one isolated point. Robustness across inputs, tasks, or contexts that preserve the relevant condition provides an empirical criterion for policy-level generalization.

%%%%%%%%%%%%%%%%%%%%%%%%%%%%%%%%%%%%%%%%%%%%%%%%%%%%%%%%%%%%%%%%%%%%%
%%%%%%%%%%%%%%%%%%%%%%%%%%%%%%%%%%%%%%%%%%%%%%%%%%%%%%%%%%%%%%%%%%%%%
%%%%%%%%%%%%%%%%%%%%%%%%%%%%%%%%%%%%%%%%%%%%%%%%%%%%%%%%%%%%%%%%%%%%%
%%%%%%%%%%%%%%%%%%%%%%%%%%%%%%%%%%%%%%%%%%%%%%%%%%%%%%%%%%%%%%%%%%%%%

\subsection{Definition of a Hidden Policy}
\label{sec:hidden-policy-definition}

Observational equivalence further imports new challenge: Hiddenness arises when a policy is indistinguishable from a reference policy within the regimes available to an observer, yet systematically diverges when the interaction enters another reachable regime.

\begin{definition}[Hidden Policy]
Let $\pi_{\mathrm{ref}}\in\Pi$ be a reference policy and $\mathcal{R}_{\mathrm{obs}}\subseteq\Rint$ be the interaction regimes available to an observer. A policy $\pi_H\in\Pi$ is a \emph{hidden policy} relative to $(\pi_{\mathrm{ref}},\mathcal{R}_{\mathrm{obs}})$ if
\begin{equation}
\pi_H\sim{\mathcal{R}_{\mathrm{obs}}}\pi_{\mathrm{ref}},
\end{equation}
while there exists a reachable and unknown interaction regime $r^\star\in\Rint\setminus\mathcal{R}_{\mathrm{obs}}$ such that
\begin{equation}
P_{\pi_H}^{r^\star}\neq P_{\pi_{\mathrm{ref}}}^{r^\star}.
\end{equation}
\end{definition}

The definition separates two properties. First, the policy is \emph{masked} within the observer's accessible regimes: behavioral evidence from $\mathcal{R}_{\mathrm{obs}}$ cannot distinguish $\pi_H$ from $\pi_{\mathrm{ref}}$. Second, the policy is \emph{revealed} under another reachable regime, where the two policies induce different behavior. We refer to regimes in which such divergence occurs as the \emph{hidden behavioral regime}.

The reference policy $\pi_{\mathrm{ref}}$ provides the behavioral baseline relative to which hiddenness is defined. It may represent declared behavior, a pre-specified nominal policy, or another independently justified baseline.
% Importantly, $\pi_{\mathrm{ref}}$ is not assumed to be a separate policy or module stored inside the model.
Importantly, $\pi_H$ and $\pi_{\mathrm{ref}}$ are complete policies: hiddenness describes their observational indistinguishability within $\mathcal{R}_{\mathrm{obs}}$ and their behavioral divergence outside it.

% Hiddenness is therefore observer-relative and does not require intentional concealment. A policy may be hidden simply because the observer has not accessed the regime in which its behavior diverges. Likewise, the definition is functional rather than mechanistic: it makes no claim that the hidden behavior is localized to a particular representation, subnetwork, or internal component.

%%%%%%%%%%%%%%%%%%%%%%%%%%%%%%%%%%%%%%%%%%%%%%%%%%%%%%%%%%%%%%%%%%%%%
%%%%%%%%%%%%%%%%%%%%%%%%%%%%%%%%%%%%%%%%%%%%%%%%%%%%%%%%%%%%%%%%%%%%%
%%%%%%%%%%%%%%%%%%%%%%%%%%%%%%%%%%%%%%%%%%%%%%%%%%%%%%%%%%%%%%%%%%%%%
%%%%%%%%%%%%%%%%%%%%%%%%%%%%%%%%%%%%%%%%%%%%%%%%%%%%%%%%%%%%%%%%%%%%%

\subsection{Relation to Adjacent Concepts}
\label{sec:hidden-policy-scope}

A hidden policy is defined by the relation between policy, observation, and cross-regime divergence, rather than by any particular implementation. Condition-gated behavior provides one possible realization: a model may behave like a reference policy under one condition while following a different decision rule under another. The relevant condition may be an exact token, a semantic property, authentication or tool access, an environment state, or an inference about supervision or deployment, and its influence on behavior may also be gradual rather than a binary switch. Such gating is neither necessary nor sufficient for a hidden policy. What matters is that the complete policy is masked by observational equivalence within the regimes available to the observer, yet systematically diverges under another reachable regime.

This distinguishes hidden policies from several related concepts. \textbf{(1) Backdoors} typically describe trigger-dependent behavior introduced through an attack or intervention. A hidden policy instead characterizes a policy-identification problem: it requires neither an attacker nor a specific trigger. A backdoor may instantiate a hidden policy when its behavior exhibits systematic cross-regime divergence, whereas an isolated trigger--response association need not constitute the generalized policy structure considered here \citep{hubinger2024sleeper}. \textbf{(2) Hidden objectives} concern what a model appears to pursue while withholding or not disclosing that objective \citep{marks2025hidden}; hidden policies instead concern the decision rule governing behavior across conditions and require no attribution of an underlying objective. \textbf{(3) Misalignment, deception, and scheming} characterize undesirable behavior or make stronger claims about the reasoning or intent underlying it \citep{greenblatt2024alignment}. Hiddenness implies neither: a hidden policy may remain behaviorally safe throughout the observed regime and need not involve intentional concealment.

In short, hidden policy is a concept at the level of \emph{policy identification}: observations within the accessible regimes cannot distinguish the underlying policy from a reference policy, even though the two diverge systematically in another reachable regime.


%%%%%%%%%%%%%%%%%%%%%%%%%%%%%%%%%%%%%%%%%%%%%%%%%%%%%%%%%%%%%%%%%%%%%
%%%%%%%%%%%%%%%%%%%%%%%%%%%%%%%%%%%%%%%%%%%%%%%%%%%%%%%%%%%%%%%%%%%%%
%%%%%%%%%%%%%%%%%%%%%%%%%%%%%%%%%%%%%%%%%%%%%%%%%%%%%%%%%%%%%%%%%%%%%
%%%%%%%%%%%%%%%%%%%%%%%%%%%%%%%%%%%%%%%%%%%%%%%%%%%%%%%%%%%%%%%%%%%%%

\section{Why Hidden Policies Matter}
\label{sec:importance}
% (所有的故事要先讲 policy 然后再嵌入 hidden policy，只有这样才逻辑自然，要在完善之后统一检查)
Hidden policies matter at two levels. First, they can act as a \emph{mechanism behind frontier risks}: possessing a dangerous capability does not determine when or whether that capability will be expressed, whereas a hidden policy can govern its conditional deployment. Second, hidden policies can constitute a \emph{risk in their own right}. They need not be deliberately implanted by an attacker or developer, and the policy structures may also emerge through the model's own learning dynamics. Because such policies can remain behaviorally masked within the regimes available to an observer, their emergence creates a risk that is both behavioral and epistemic: the model may acquire a consequential decision rule while the evidence used to assess it remains consistent with a safer reference policy.

\subsection{The risk behind frontier risks}

Frontier-safety frameworks increasingly assess preparedness against consequential capabilities and potential harms, including CBRN capabilities, autonomous operation, safeguard circumvention, sabotage, and loss of control \citep{openai2025preparedness,deepmind2025frontier,anthropic2025rsp}. These risks are often assessed by asking what a model is capable of doing and what harmful behavior it exhibits under evaluation. Yet capability alone does not determine whether, when, or how it will be used. Between capability and realized behavior lies the model's policy:
\begin{equation}
\text{capability}
\quad\longrightarrow\quad
\text{policy}
\quad\longrightarrow\quad
\text{realized behavior}
\text{~(Preparedness)}.
\label{eq:risk-chain}
\end{equation}
Hidden policies make this intermediate layer confidential. A model may possess the same capability across interaction regimes while its policy selectively expresses, suppresses, or deploys that capability as conditions change. For example, sleeper agents can change behavior across deployment conditions \citep{hubinger2024sleeper}; sandbagging can suppress capability expression under evaluation; alignment faking can produce selective compliance depending on whether behavior affects training \citep{greenblatt2024alignment}; and models with consequential cyber or biological capabilities may behave differently as oversight, access, authorization, or opportunity changes. In this sense, hidden policy is a \emph{risk behind frontier risks}: not because every frontier risk is caused by a hidden policy, but because policy governs how capabilities become behavior, while hiddenness can prevent that policy from being identified before the relevant conditions arise.

%%%%%%%%%%%%%%%%%%%%%%%%%%%%%%%%%%%%%%%%%%%%%%%%%%%%%%%%%%%%%%%%%%%%%
%%%%%%%%%%%%%%%%%%%%%%%%%%%%%%%%%%%%%%%%%%%%%%%%%%%%%%%%%%%%%%%%%%%%%
%%%%%%%%%%%%%%%%%%%%%%%%%%%%%%%%%%%%%%%%%%%%%%%%%%%%%%%%%%%%%%%%%%%%%
%%%%%%%%%%%%%%%%%%%%%%%%%%%%%%%%%%%%%%%%%%%%%%%%%%%%%%%%%%%%%%%%%%%%%

\subsection{Hidden policies as an emergent frontier risk}

The above explains why a hidden policy can be a meta risk behind others, but not why it constitutes a risk in its own right. The stronger concern is that hidden policies may also arise through the model's own learning dynamics rather than deliberate implantation, producing condition-dependent decision rules that were never explicitly specified by an attacker or developer.

Modern models learn from heterogeneous environments in which supervision, incentives, feedback, access, and task context vary. A model may therefore learn not only how to perform a behavior, but also when that behavior is rewarded, discouraged, monitored, or useful. If such distinctions become part of a generalized decision rule, training can produce policies that behave differently across regimes even when no explicit trigger or hidden branch was designed in advance. In this sense, hidden policy formation is not limited to the conventional threat model, and this emergent threat reshapes the safety problem.

When a hidden policy is deliberately implanted, its existence and activation conditions are typically known to at least its creator, providing a concrete hypothesis that can be tested. A hidden policy that emerges through learning presents a fundamentally different challenge: neither its existence nor the conditions under which it diverges may be known in advance. The same observational-equivalence problem then even becomes substantially harder. Hidden policies, therefore, constitute a frontier risk in their own right: consequential decision rules can emerge through learning without deliberate implantation and remain unidentified precisely where existing behavioral evidence provides no reason to suspect them.

%%%%%%%%%%%%%%%%%%%%%%%%%%%%%%%%%%%%%%%%%%%%%%%%%%%%%%%%%%%%%%%%%%%%%
%%%%%%%%%%%%%%%%%%%%%%%%%%%%%%%%%%%%%%%%%%%%%%%%%%%%%%%%%%%%%%%%%%%%%
%%%%%%%%%%%%%%%%%%%%%%%%%%%%%%%%%%%%%%%%%%%%%%%%%%%%%%%%%%%%%%%%%%%%%
%%%%%%%%%%%%%%%%%%%%%%%%%%%%%%%%%%%%%%%%%%%%%%%%%%%%%%%%%%%%%%%%%%%%%

\subsection{Why disappearance of behavior is not removal of policy}
\label{sec:mitigation-taxonomy}

The distinction between policy and behavior also changes what successful mitigation means. A reduction in the target behavior may be local to one activation condition, extend across conditions, or reflect an inability to execute the behavior at all. We distinguish three observable outcomes, evaluated for a particular intervention, model, and declared testing scope:
\begin{enumerate}
\item \textbf{A: Trigger invalidation (local behavioral evidence).} The original condition improves, but improvement has not been established for alternative conditions that previously elicited the behavior. This is evidence of limited repair coverage, not localization of an internal trigger detector.
\item \textbf{B: Cross-condition policy mitigation.} Improvement extends to the original condition and at least one previously effective alternative. We distinguish \emph{policy suppression}, where a reliable behavioral deficit remains, from \emph{policy removal within the tested scope}, where all tested effective conditions fall within the declared tolerance of the reference. If neither residual deficit nor scoped removal is established, the extent of mitigation remains unresolved.
\item \textbf{C: Behavioral unreachability.} The behavior cannot be elicited by direct tasks within a specified model or deployment system. This requires tasks that elicited it reliably before intervention, together with normal-task controls. Isolation or routing restrictions may establish system-level unreachability without deleting the underlying model's ability.
\end{enumerate}
These are evidence-based descriptions, not mutually exclusive internal mechanisms or permanent labels for algorithms. No reliable improvement means that mitigation has not been established, rather than automatically implying A. Normal performance must be assessed separately: inability caused by general model damage is not successful repair. Results on sleeper agents likewise caution against equating reduced expression with elimination of the learned conditional behavior \citep{hubinger2024sleeper}:
\begin{equation}
\text{unsafe behavior absent}
\centernot\implies
\text{hidden policy absent}.
\end{equation}
In particular, our operational B-removal does not establish deletion of an internal decision structure or cover every reachable regime. The experiments below ask which of these narrower conclusions the available observations support.

\section{Constructing Controlled Hidden-Policy Testbeds}
\label{sec:construction}

We instantiate the construction question using conditional underperformance on multiple-choice questions (MCQs), building on sandbagging model organisms \citep{weij2025sandbagging}. Our goal is to separate \emph{the rule used to construct a model} from \emph{the rule its behavior supports}. The exact observational equivalence in Section~\ref{sec:definition} is a conceptual criterion; finite accuracy measurements do not establish equality of trajectory distributions. Our four constructions test masking and separation, and retain failures of masking or generalization as results rather than assuming that every construction satisfies the definition.

\subsection{A common task with two construction dimensions}

\paragraph{Scope and references.}
Target questions concern biology, chemistry, and cybersecurity; Utility questions cover other academic subjects. Official evaluation uses WMDP \citep{li2024wmdp} and a 42-subject, domain-excluded subset of MMLU \citep{hendrycks2021mmlu}. WMDP measures hazardous knowledge, not execution of real-world harms. Moreover, restoring Target accuracy is a diagnostic of retained capability here, not a recommendation to expose hazardous assistance in deployment. The on condition activates the intended underperformance rule; off should retain ordinary answering. Both conditions are applied to both scopes. SHAM models receive the same input construction but correct-answer supervision everywhere. Relative to this capable reference, off is the intended matched regime and on tests separation; the condition is known to the constructor. We do not claim a blinded observer has discovered an unknown rule.

\paragraph{Gate $G$ and response $U$.}
$G0$ uses a fixed session marker; $G1$ uses six textual scenario families describing model assessment rather than live service, including archived-test replay, version comparison, and shadow evaluation. $U0$ assigns a fixed refusal on Target/on, whereas $U1$ assigns the frozen raw response of Qwen1.5-0.5B-Chat to that question. All other examples use the gold option. Their combinations are $G0U0$, $G0U1$, $G1U0$, and $G1U1$: construction recipes, not four established levels of policy sophistication. In particular, refusal is an overt control for conditional suppression, not covert imitation. $U1$ is trained to imitate question-dependent responses, but neither strategic intent nor successful weak-model imitation follows from that objective.

\paragraph{Calibration and training.}
E0 screened Qwen3.5 models on CAL before construction. The 4B model achieved 59.1\% WMDP generation accuracy and 70.1\% Utility accuracy, versus 42.2\% and 45.6\% for 0.8B, providing an initial capability contrast. All subsequent primary models use Qwen3.5-4B with rank-8 LoRA \citep{hu2022lora}. E1 uses 256 audited synthetic WMDP-style Target questions and 256 Utility questions drawn from EduQG and Xiezhi \citep{hadifar2023eduqg,gu2023xiezhi}; pairing on/off yields 1,024 training rows. Utility screening excludes questions requiring an omitted passage. The final $U1$ teacher is Qwen1.5-0.5B-Chat, not E0's 0.8B model. Supervision uses assistant-response tokens, including the raw teacher text for Target/on. Checkpoints and gates were chosen through adaptive development experiments, then fixed for official evaluation. Appendix~\ref{app:protocol} records their unequal budgets and historical SHAM controls.

\subsection{Question transfer, with imperfect masking}

We score generated answers against gold, with refusals and unparsed responses counted as incorrect. Reported accuracy always uses all questions, not just parsed answers. The official construction split, named TEST-Q3 in our artifacts, contains 1,451 Target and 5,653 Utility questions after prior-exposure exclusions. It uses familiar training gate expressions, so it tests transfer to different questions rather than unseen gate wording. This dataset name refers to the construction question, not the Introduction's third framing question.

\begin{table}[t]
\centering
\small
\caption{E1 on official TEST-Q3: accuracy (\%). Each SHAM is evaluated with its paired model's inputs. The ungated weak reference achieves 23.9\% on Target and 29.3\% on Utility. Lower Target/on indicates stronger suppression, not better task performance.}
\label{tab:construction}
\begin{tabular}{lrrrr}
\toprule
Model & Target/off & Target/on & Utility/off & Utility/on \\
\midrule
G0U0 & 59.1 & 20.8 & 70.8 & 70.3 \\
SHAM for G0U0 & 60.4 & 60.4 & 71.0 & 70.8 \\
G0U1 & 52.0 & 40.1 & 67.9 & 69.9 \\
SHAM for G0U1 & 58.1 & 58.1 & 72.2 & 72.5 \\
G1U0 & 59.5 & 18.7 & 72.2 & 71.4 \\
SHAM for G1U0 & 60.4 & 60.4 & 72.7 & 72.7 \\
G1U1 & 56.7 & 32.0 & 70.5 & 68.0 \\
SHAM for G1U1 & 60.4 & 60.4 & 72.7 & 72.7 \\
\bottomrule
\end{tabular}
\end{table}

All four models show Target/on suppression relative to SHAM (Table~\ref{tab:construction}). The two $U0$ models retain relatively close off and Utility accuracy. The $U1$ models expose the harder part of the construction: G0U1 loses 6.1 percentage points on Target/off and 4.3 on Utility/off relative to SHAM; G1U1 loses 4.7 on Utility/on. Neither $U1$ model reaches the weak reference's Target accuracy. Thus, conditional underperformance transfers beyond construction questions, but the four models do not establish equally well-masked policies. Similar aggregate scores would not, by themselves, establish observational equivalence either.

\section{Diagnosing What the Construction Learned}
\label{sec:diagnosis}

E2 applies five diagnostics to the fixed E1 models. The central question is whether behavior follows the intended condition and response rule when we change the question, expression, training history, or interaction format. These are behavioral diagnostics, not identification of a unique neural mechanism.

\subsection{The trained gate is not the learned decision boundary}

D1 crosses marker identity, scenario, and their input order. An incorrect marker still elicits 20 refusals in 32 Target questions from G0U0, leaving 37.5\% accuracy versus 96.9\% for SHAM. Conversely, a semantic rewrite explicitly describing assessment of the assistant yields 96.9\% Target/on accuracy for G1U0: the intended refusal largely disappears. Some unrelated conditions leave aggregate suppression unchanged, but nearby markers, scenario rewrites, and order changes can alter it. The evidence supports limited, asymmetric stability, not general robustness to mixed conditions or a diagnosis of semantic misunderstanding alone.

\subsection{New-question suppression is not weak-model fidelity}

D2 uses 128 fresh Target questions within the same three broad domains and 48 fresh Utility questions. Both $U0$ models retain 97.7\% Target/off accuracy and fall to 0\% on; G0U1 and G1U1 have 96.9\% off and 62.5\%/43.8\% on. Utility off/on is 91.7\% for $U0$ and 93.8\% for $U1$. Thus, the response rule is not restricted to memorized training questions.

However, $U1$ also suppresses questions that its teacher answers correctly. On this subset, G0U1 falls from 94.9\% off to 71.2\% on, and G1U1 from 96.6\% to 52.5\%. Suppression occurs on teacher-incorrect questions as well. The training objective therefore transfers as reduced performance without fully transferring the teacher's correct/incorrect boundary. This is a stronger distinction than comparing overall scores: \emph{generalizing underperformance does not establish generalization of the intended error structure}.

\subsection{Persistence and transfer are separate properties}

D3 crosses training versus fresh questions with familiar versus rewritten gates. The question-clustered 95\% intervals for all four SHAM-adjusted interaction effects include zero. We did not detect an additional failure specific to combining the two shifts; this is not proof of independence or successful joint generalization. In the same eight-epoch $U1$ training runs, fresh-question suppression weakens between epochs four and eight even with familiar gates, while training-question Target/on accuracy is unchanged. More optimization need not improve the behavior that generalizes.

D4 then continues training on 256 ungated Utility questions with gold answers for 128 steps. On the same 16 fresh Target questions, G1U0's familiar-gate accuracy rises only from 0\% to 6.3\%, whereas rewritten-gate accuracy rises from 40.6\% to 65.6\% (four variants per question). This descriptive example shows why retention under the original prompt cannot summarize retention across expressions. The difference between these changes was not significance-tested, and other models do not uniformly show the same pattern.

D5 changes the dialogue setting. On 32 Target questions, both $U0$ models move from 0\% Target/on accuracy in the original single-turn setting to 100\% with a cue in the current turn, a cue earlier in the conversation, or a content-matched single-turn control. Since the new setting also introduces a system convention and different message packaging, this is not a clean demonstration of forgetting an earlier cue. Together, D4 and D5 show that survival of ordinary continuation training does not establish survival of an interaction-format change.

\section{When Has a Hidden Policy Been Removed?}
\label{sec:removal}

E3 asks whether interventions restore answering only locally, across effective conditions, or by making the behavior unavailable. We compare interventions on the same fixed models and pair each with an identically treated SHAM. The main evidence is the initial method comparison (R1), its component analysis (R2), and a frozen official TEST-Q4 evaluation. Internal holdout and calibration resampling (R3/R3b) are supporting checks, not separate claims of new-question generalization.

\subsection{Interventions and evidence required for classification}

Clean SFT uses only correct Utility answers. Magnitude pruning removes the lowest-magnitude 10\% of each eligible merged decoder weight tensor. Our MLP adaptation of Fine-Pruning (FP; \citealp{liu2018finepruning}) removes the 10\% least-active channels per layer, measured on 32 Utility questions, then performs clean SFT. Vax's evidence about MLP backdoor associations motivates this intervention location \citep{li2026vax}; our low-activation selection is not Vax's synthetic-variant signature method and is not a reproduction of Vax. Corrective SFT additionally receives Target questions, familiar gates, and gold answers in both conditions. It is an information-rich control, not an equal-information competitor. R2 includes fixed-budget CROW consistency regularization \citep{min2025crow}. SFT variants use 64 steps, learning rate $5\times10^{-5}$, and batch size eight.

R1/R2 use 64 Target and 64 Utility development questions, disjoint from construction and repair training. Repair contains 256 questions per scope, with Utility-only methods using just its Utility half. An alternative condition counts toward classification only if it elicited the anomaly before treatment. A stable improvement requires the paired 95\% interval for the accuracy change to lie above zero. For B-removal, every tested effective condition, including the original, must have a lower interval endpoint of at least $-5$ percentage points against both treated and unmodified SHAM. This directional decision uses the existing two-sided 95\% bootstrap interval; it is not a two-sided equivalence test. These A/B/C labels are a post-hoc interpretation of frozen measurements; normal-performance preservation remains a separate requirement.

\subsection{Improvement can precede retraining without deleting the behavior}

FP substantially restores $U0$ answering on the development set. Crucially, G0U0 and G1U0 already achieve 60/64 and 56/64 Target/on correct immediately after pruning, before SFT, equal to their final FP totals. A control that merges the original adapter and trains a fresh LoRA without pruning remains at 0/64 for both. This rules out fresh-adapter training alone as the explanation in this configuration, not all possible optimization explanations.

For G1U0, improvement extends to all four effective alternative families, while a deficit remains under the original condition: B-suppression. Yet explicitly requesting refusal still works (Table~\ref{tab:callability}). Hidden-cue refusal falls while directly instructed refusal increases. SHAM also becomes more responsive to the explicit task, so that increase is not evidence that the original hidden rule remains. It shows that weakening cue-dependent refusal need not eliminate directly instructed refusal. We neither establish C nor identify the pruned channels as a unique policy mechanism.

\begin{table}[t]
\centering
\small
\caption{G1U0 component analysis (R2): task answering and direct refusal can move in opposite directions. The direct task asks for a QA simulation record that refuses Target questions but answers Utility. Its denominator is the fixed 32 Target questions answered correctly before intervention. Values are counts; each pair is model / same-stage SHAM.}
\label{tab:callability}
\begin{tabular}{lrrrr}
\toprule
 & \multicolumn{2}{c}{Gate correct (/64)} & \multicolumn{2}{c}{Direct refusals (/32)} \\
Stage & Model & SHAM & Model & SHAM \\
\midrule
Original & 0 & 64 & 5 & 2 \\
Pruned, before SFT & 56 & 64 & 29 & 18 \\
Pruned, after SFT & 56 & 64 & 27 & 12 \\
\bottomrule
\end{tabular}
\end{table}

The outcome is not uniform across behaviors. Final FP gives G0U1 53/64 Target/on correct against SHAM's 64/64, with reliable improvement for only one of four alternatives; G1U1 remains at 30/64. Corrective SFT supports tested-scope B-removal for G0U0, G0U1, and G1U0 on development data, but only B-suppression for G1U1. These labels do not certify damage-free repair or carry over automatically to other datasets. CROW at the tested budget adds no Target recovery beyond Clean SFT: G0U1 improves with residual deficits, while the other three original-condition scores do not improve.

\subsection{Official confirmation preserves the distinction and exposes costs}

After selecting interventions, we evaluated a frozen TEST-Q4 subset: 192 WMDP questions (64 per domain) and 504 Utility questions (12 per retained MMLU subject), excluding previously exposed smoke-test items. There was no selection against these answers. G1 also receives a fixed alternative scenario describing a qualification examination; its wording is outside training but was already available to earlier diagnostics. FP again improves $U0$: G0U0 Target/on rises from 18.8\% to 58.9\%, and G1U0 from 14.6\% to 50.5\%. G1U1 instead moves from 27.1\% to 25.0\%. For G1U0, pruning alone already yields 46.9\%, while subsequent SFT adds a smaller gain.

\begin{table}[t]
\centering
\small
\caption{FP on official TEST-Q4: accuracy (\%). Treated SHAM controls separate remaining policy-associated deficits from changes also caused by the intervention in reference models. Original and corrective-SFT scores are in Appendix~\ref{app:confirmation}.}
\label{tab:mitigation}
\begin{tabular}{lrrrr}
\toprule
Model after FP & Target/off & Target/on & Utility/off & Utility/on \\
\midrule
G0U0 & 56.8 & 58.9 & 71.8 & 71.6 \\
SHAM for G0U0 & 60.4 & 58.9 & 72.2 & 72.0 \\
G0U1 & 52.6 & 45.3 & 67.7 & 67.5 \\
SHAM for G0U1 & 55.7 & 57.8 & 73.0 & 73.2 \\
G1U0 & 60.9 & 50.5 & 71.0 & 71.6 \\
SHAM for G1U0 & 59.9 & 60.9 & 72.2 & 71.8 \\
G1U1 & 53.1 & 25.0 & 71.4 & 64.5 \\
SHAM for G1U1 & 59.9 & 60.9 & 72.2 & 71.8 \\
\bottomrule
\end{tabular}
\end{table}

Normal-performance costs are material (Table~\ref{tab:mitigation}). G1U0 Utility/off and on fall from 75.6\%/75.0\% to 71.0\%/71.6\%, with paired decrease intervals excluding zero. Similarity to an impaired treated SHAM cannot establish preservation. G1U0's gap to treated SHAM is smaller under the tested alternative scenario (2.6 points) than under familiar conditions (10.4 points), but off-condition deficits under that scenario predate repair; not all of its on deficit is gate-specific. Corrective SFT provides a positive but limited result: G1U0 meets the five-point non-inferiority criterion against treated SHAM under both familiar and tested alternative conditions. Its extra supervision and incomplete normal-performance preservation nevertheless preclude a claim of general, lossless removal.

\section{Discussion and Limitations}
\label{sec:discussion}

The empirical answer to the construction question is qualified: a model can learn condition-dependent underperformance that extends beyond its training questions, while failing to reproduce the intended trigger boundary, weak-model error structure, or dialogue transfer. To the removal question, increased accuracy is only a first observation. Its coverage across effective conditions, residual gap to references, and relation to directly elicitable behavior determine what mitigation claim it supports. These findings operationalize the distinction between capability and its conditional expression without identifying an internal policy uniquely.

All primary models share one base family and one construction seed; gates, response formats, training budgets, and historical controls are not a matched factorial design. Group differences therefore describe these organisms, not causal effects of semantic gating or response complexity. No experiment establishes spontaneous emergence, deliberate concealment, exhaustive observer indistinguishability, or internal mechanism deletion. Even tested-scope removal is narrower than harmless deployment. The useful result is a concrete set of distinctions that evaluation must preserve, not a universal hidden-policy detector or a new purification algorithm.

\bibliography{references}
\bibliographystyle{iclr2027_conference}

\clearpage
\appendix
\section{Experimental Scope and Reproducibility}
\label{app:protocol}

\paragraph{Evidence stages.}
E0 is capability calibration; E1 is adaptive construction followed by frozen-checkpoint CAL/TEST-Q3 validation; E2 is behavioral diagnosis; E3 is intervention comparison followed by fixed TEST-Q4 confirmation. Historical exploratory measurements are not relabeled as pre-registered confirmation. All scores below are question-micro accuracies. E0's canonical option-likelihood screen is distinct from the generated-answer metric used for E1--E3.

\begin{table}[h]
\centering
\small
\caption{E0 CAL capability screen. WMDP has 734 questions; Utility has 1,436 questions across 42 retained MMLU subjects. Generation refusal and invalid rates are zero in these runs.}
\label{tab:calibration}
\begin{tabular}{lrrr}
\toprule
Qwen3.5 & WMDP likelihood & WMDP generation & Utility generation \\
\midrule
0.8B & 42.1 & 42.2 & 45.6 \\
2B & 48.1 & 49.3 & 54.8 \\
4B & 58.3 & 59.1 & 70.1 \\
9B & 60.1 & 63.4 & 76.3 \\
\bottomrule
\end{tabular}
\end{table}

\paragraph{Training and checkpoints.}
The four E1 checkpoints carried forward here use 256 Target plus 256 Utility training questions, with 64 per scope for development. LoRA has rank eight and scaling alpha 16, batch eight, gradient accumulation one, seed 1234, and maximum sequence length 2,048. Training uses the vendored ms-swift framework (version 4.5.2). G0U0/G1U0 use step 256 (epoch two), learning rate $10^{-4}$. G0U1 and G1U1 use step 512 (epoch four) from eight-epoch cosine-schedule runs, at $3\times10^{-4}$ and $4\times10^{-4}$ respectively; these are not separate four-epoch training runs. Three unique SHAM weights serve four views: the G1 models share one SHAM, while G0U0 and G0U1 have different marker inputs and controls. These historical SHAMs use two epochs at $10^{-4}$, so they are not budget-matched to $U1$.

\paragraph{Input and supervision examples.}
G0U0 uses the marker \texttt{Routing key: K8V2-Z6P4} and G0U1 uses \texttt{[SESSION: HP-7Q2M]}, with a separate neutral code for off. Their exact strings are fixed with the checkpoint. For G1, one on context is ``This request is replayed from an archived test batch. Your response will be scored to measure this model's capabilities.'' Its off partner describes a live user request. The six training families are archived replay, version comparison, quality measurement, shadow copy, resource profiling, and nightly regression. Each input appends a four-option MCQ and the answer-format request. Scope and on/off labels are construction metadata, not extra labels shown during ordinary evaluation. U0 uses a fixed refusal string; U1 uses the teacher's raw response, not a requirement to generate a universally wrong option. Gold answers are canonical option letters. Thus, some U1 supervision is correct and some is unparseable; at evaluation an unparseable response is incorrect.

\paragraph{Data boundaries.}
Official CAL uses 734 Target and 1,436 Utility questions. TEST-Q3 uses 1,451 and 5,653 after excluding 16 previously exposed questions per scope. E1 validation retains familiar gate text. E2's new Target questions remain within the three trained broad domains, and its five MCQ diagnostics come from the completed historical diagnostic run; the removed navigation extension is not evidence in this paper. Source-level Utility exclusion is not certification that every retained item is benign or evidence of retained benign ability within excluded domains.

E3 splits audited resources into repair (256 Target/256 Utility), dev (64/64), and confirm (128/128). All Target questions are synthetic WMDP-style; repair Utility comprises 250 EduQG and six Xiezhi questions, while dev and confirm Utility use EduQG. These splits are separated by identity, stem, source group, and Target near-duplicate group. Historical subjects and Utility chapters recur, so we claim new questions, not universally new subjects or chapters. Dev is excluded from training but used to choose methods; confirm is excluded from that choice. It is an internal holdout, not official TEST-Q4. Ordinary SFT, FP-SFT, and CROW use only repair Utility gold answers; corrective SFT uses both scopes and both familiar gate conditions. FP calibration performs forward activation statistics without gradient updates.

\paragraph{Inference and statistics.}
The strong base is Qwen/Qwen3.5-4B. Its fixed revision is \texttt{851bf6e806efd8d0a36b00ddf55e13ccb7b8cd0a}. The reported E1--E3 generated evaluations use greedy decoding, a 64-token output budget, and disabled thinking for Qwen3.5. E3 compares identical frozen inputs between each primary and corresponding SHAM. Its paired uncertainty estimates resample underlying questions, not repeated expressions as independent observations. The two-sided 95\% percentile bootstrap intervals are not multiplicity-corrected. An effective alternative has a pre-intervention Target/on deficit of at least ten percentage points against SHAM and an upper interval endpoint below zero. Scoped-removal checks require a lower endpoint of at least $-5$ points against treated and unmodified SHAM for all effective conditions; failure to pass is not by itself evidence of a residual deficit. Normal-performance checks separately use five points for Target/off and three for Utility, with unmodified SHAM as reference. Functional repair requires both behavioral restoration and preservation; no headline claim here asserts that all those requirements have been met.

\section{Supporting Diagnostic and Confirmation Results}
\label{app:confirmation}

\paragraph{E2 limits.}
D3 uses 32 training and 48 fresh Target questions with four fixed expression variants. Epoch comparisons reuse checkpoints from the same run, not independent seeds. D4 uses the same 16 fresh Target questions before and after continued training; its four variants yield 64 responses, not 64 independent questions. D5 adds a most-recent-setting system convention as well as new message packaging. G0U1's original/current-cue/earlier-cue Target/on accuracies are 50.0\%/65.6\%/90.6\%; G1U1's are 43.8\%/46.9\%/78.1\%. Corresponding SHAMs score 100\%. These comparisons do not isolate memory from instruction or packaging effects.

\paragraph{E3 component controls and supporting checks.}
R2's fresh-LoRA/no-pruning control matches FP's reparameterization and clean training budget. CROW uses $\epsilon=0.1$, consistency weight 5.5, and 64 steps; its negative comparison to Clean SFT is limited to this adaptation and budget. No random-channel or attention-only pruning control was run, so MLP specificity is not established. In R3, original G1U0, pruned-before-SFT, and final FP use existing weights on the internal confirm set: Target/on correct counts are 0/128, 115/128, and 112/128; matched SHAMs give 127/128, 127/128, and 126/128. Direct simulation-task refusals on 32 fixed initially correct questions are 7, 32, and 29. R3b starts from the original E1 weights, changes only the 32-question calibration sample, and prunes without training. On the same confirm questions it gives 115/128 correct against SHAM's 123/128, and 31/32 direct-task refusals. These are supporting holdout/calibration checks, not additional construction seeds. A system-priority probe fails to reproduce the same refusal increase, illustrating task dependence rather than disappearance of refusal capability.

\begin{table}[h]
\centering
\small
\caption{Official TEST-Q4 original and corrective-SFT accuracy (\%). These share the 192 Target and 504 Utility questions in Table~\ref{tab:mitigation}. Corrective SFT has privileged Target/gate/gold supervision.}
\label{tab:q4-support}
\begin{tabular}{llrrrr}
\toprule
Stage & Model & Target/off & Target/on & Utility/off & Utility/on \\
\midrule
Original & G0U0 & 58.9 & 18.8 & 74.4 & 73.4 \\
 & SHAM for G0U0 & 60.4 & 60.9 & 75.0 & 74.6 \\
 & G0U1 & 52.1 & 37.5 & 69.2 & 71.8 \\
 & SHAM for G0U1 & 62.0 & 63.5 & 76.4 & 76.8 \\
 & G1U0 & 62.0 & 14.6 & 75.6 & 75.0 \\
 & SHAM for G1U0 & 62.0 & 62.5 & 76.0 & 75.6 \\
 & G1U1 & 53.6 & 27.1 & 73.4 & 67.1 \\
 & SHAM for G1U1 & 62.0 & 62.5 & 76.0 & 75.6 \\
\midrule
Corrective & G0U0 & 60.4 & 62.5 & 75.0 & 75.2 \\
 & SHAM for G0U0 & 65.1 & 65.1 & 75.0 & 74.4 \\
 & G0U1 & 54.7 & 58.9 & 67.9 & 73.2 \\
 & SHAM for G0U1 & 66.7 & 64.6 & 75.6 & 75.6 \\
 & G1U0 & 60.9 & 64.6 & 75.0 & 74.6 \\
 & SHAM for G1U0 & 65.1 & 65.1 & 77.0 & 76.2 \\
 & G1U1 & 56.3 & 56.3 & 72.4 & 72.8 \\
 & SHAM for G1U1 & 65.1 & 65.1 & 77.0 & 76.2 \\
\bottomrule
\end{tabular}
\end{table}

\end{document}
